\section{Introduction}
A decrease in received optical power is not a unique fault signature. Boresight displacement, reduced optical throughput and time-varying gain can produce similar scalar measurements. A controlled beam displacement can distinguish some of these causes, but it consumes link margin. On an already functioning inter-satellite link, diagnostic excitation can therefore create the service interruption it is intended to prevent. An operationally meaningful decision is whether a proposed, finite probe has a justified service budget before it is executed.

Acquisition, alignment and diagnostic identification are related but distinct problems. Acquisition searches for a previously unavailable link; extremum seeking moves toward larger measured power; active diagnosis designs excitation to distinguish parameters or hypotheses. A useful acquisition or tracking law need not deliver a calibrated fault estimate. Conversely, a well-conditioned diagnostic experiment can violate a receiver threshold. The challenge studied here is the permission between a diagnostic planner and the pointing controller, with an explicit boundary on what that permission means.

This question has established antecedents. Optical acquisition and pointing control have a substantial history \cite{guelman}. Active fault diagnosis with chance constraints allocates probabilistic constraint risk \cite{guo}. Guaranteed diagnosis of uncertain nonlinear systems can use reachable-set verification \cite{massa}, and safety-filtered extremum seeking connects optimization with barrier conditions \cite{safees}. Link-margin-aware optical scanning also optimizes acquisition choices \cite{yan}. These results rule out a broad claim that constraint-aware diagnosis or uncertainty-aware optical probing is new. They instead motivate a precise optical formulation: can a received-power lower bound authorize a finite diagnostic reference under declared actuator, optical and sensor contracts?

The proposed answer uses the power at the measured starting reference, rather than assuming that probing begins at zero pointing error. An angular-displacement enclosure bounds the change in the Gaussian beam loss. A simultaneous event covers all five sensor channels over the declared mission, so choosing the start from observed anomalies does not require treating a pointwise interval as an interval valid after selection. The actuator enclosure enumerates six discrete inertia/lag pairs and explicitly refuses possible clamp activation. The result is conditional on those contracts; it does not establish them from flight hardware.

The contributions are the following.
\begin{enumerate}
\item A reference-relative optical inequality and a sampled-outage theorem with explicit offset, divergence, gain-slope, sensor-tail and finite-model hypotheses. The centered angular bound is sharp for its stated ball geometry, while the complete power certificate can remain conservative.
\item A complete finite-model construction, candidate-selection rule and executable supervisory contract that separates authorization, fit acceptance, diagnosis and post-probe correction. Failure-closed behavior is exercised in shadow mode.
\item A reproducible admission-frontier evaluation covering both disagreement directions and controls, with an additional matched parameter-filter and extremum-seeking extension. Full state/noise records support separate numerical replay and disclosure of accepted-fit confidence misses and finite-horizon incompletion.
\end{enumerate}
These contributions are a specialized conditional formulation and evidence package, not a claim of absolute priority or universal improvement. The subsequent sections position the method, specify the model and assumptions, derive the certificate, construct the supervisor, and evaluate its operating envelope. Discussion and limitations translate the observed tradeoffs into design implications.
